\section{Collective arithmetic, occurrence geometry and notation}\label{sec:notation}
\subsection{The collective map}
\begin{definition}[collective map]\label{def:collective}
Let $x=a/b>0$ be a rational number in lowest terms. Its \emph{occurrence set} is
\[
 D(x)=\{\,d\in\Z:\ d\ge2,\ d^2\mid a,\ d\equiv1 \pmod b\,\}.
\]
The state $x$ is \emph{realised} if $D(x)\ne\emptyset$. For a realised state put
\[
 L(x)=\operatorname{lcm}D(x),\qquad G(x)=\gcd\{\,d-1:\ d\in D(x)\,\},\qquad C(x)=x+\frac{L(x)^2}{G(x)} .
\]
The map $C$ is the \emph{collective map} of the rational dynamics FIELD-I~\cite{FoundII}, in the form used by the programme (frozen releases~\cite{RefColl,HD1}). An \emph{orbit} is a sequence $x_0,x_1=C(x_0),\dots$ as long as the states are realised; $D_n=D(x_n)$, $L_n$, $G_n$ are written with the index. For an integer $x=m$ ($b=1$) the set $D(m)$ is the set of divisors $\ge2$ of the \emph{square-root content} $\sigma(m)=\prod_pp^{\lfloor v_p(m)/2\rfloor}$.
\end{definition}

\begin{lemma}[basic facts]\label{lem:basic}
Let $x=a/b$ be realised, $L=L(x)$, $G=G(x)$. Then $b\mid G$, every $d\in D(x)$ satisfies $d\equiv1\pmod G$, and $\gcd(L,G)=1$. Writing $G=bh$ one has $x=ah/G$ and $C(x)=(ah+L^2)/G$.
\end{lemma}
\begin{proof}
Every $d\in D$ is $\equiv1\pmod b$, so $b\mid d-1$ for all $d$ and $b\mid G$. Each $d-1$ is divisible by $G$ by definition of $G$, so $d\equiv1\pmod G$ and $\gcd(d,G)=1$; a prime dividing $L$ divides some $d$, hence does not divide $G$. Finally $a/b=ah/G$ and $L^2/G$ has the common denominator $G$. \qedhere
\end{proof}

\subsection{Exponent box, lattice, Graver set, move complex}
Write $a=\prod_pp^{v_p}$ and $q_p=\lfloor v_p/2\rfloor$. The \emph{active primes} are those with $q_p\ge1$; the \emph{box} is $B=\prod_{p\text{ active}}[0,q_p]\subset\Z^{k}$, $k$ the number of active primes, and for $\alpha\in B$ we write $p^\alpha=\prod p^{\alpha_p}$. Since $d^2\mid a\iff$ every exponent of $d$ is at most $q_p$, the map $\alpha\mapsto p^\alpha$ identifies the \emph{fibre} $F(x)=\{\alpha\in B\setminus\{0\}:p^\alpha\equiv1 \pmod b\}$ with $D(x)$; the origin $0\in B$ corresponds to $d=1$ and is never a vertex.

\begin{definition}[lattice, region-Graver set]\label{def:graver}
The \emph{relation lattice} of $x$ is $\Lambda_b=\{g\in\Z^k:\ \prod_pp^{g_p}\equiv1\pmod b\}$ (the primes are prime to $b$ because $a/b$ is reduced); $\Lambda_1=\Z^k$. For $g,h\in\Z^k$ write $h\sqsubseteq g$ if $h_ig_i\ge0$ and $|h_i|\le|g_i|$ for all $i$ (conformal order). The \emph{region-Graver set} $\Gr(x)$ consists of the $\sqsubseteq$-minimal non-zero elements of $\Lambda_b\cap[-q,q]$ (where $q=(q_p)$); it is a finite set closed under $g\mapsto-g$, a lattice-intrinsic object in the sense of Graver~\cite{Graver1975,Sturmfels1996}.
\end{definition}

\begin{definition}[move complexes]\label{def:Kmove}
A \emph{cell} $P(u;g_1,\dots,g_j)$ is the parallelotope $\{u+\sum_i\varepsilon_ig_i:\varepsilon\in\{0,1\}^j\}$ with $g_i\in\Gr(x)$ lexicographically positive and linearly independent; it is \emph{admissible} if all its $2^j$ vertices lie in $F(x)$. The \emph{move complex} $\Kmove(x)$ is the cubical cell complex of all admissible cells, with the usual cubical boundary $\partial P(u;g_1,\dots,g_j)=\sum_i(-1)^{i-1}\bigl(P(u+g_i;\hat g_i)-P(u;\hat g_i)\bigr)$ (see~\cite{Kozlov2008} for cubical chain complexes). The \emph{axis complex} $\Kaxis(x)$ is defined in the same way with the generators restricted to the unit vectors. Homology is taken with coefficients in a field; the Betti numbers recorded in this manuscript were computed over $\F_2$, and over $\F_p$ for $p=2^{31}-1$ in all cases where the two were compared they agree.
\end{definition}
The complexes $\Kmove$ and $\Kaxis$ are \emph{not} identified with the full realised complex of Manuscript~I nor with any history complex of Manuscript~II (\cref{prop:A4T7}).

\begin{remark}[Betti numbers as exact data]\label{rem:exact}
Every Betti number in a table or figure of this manuscript is a rank computed with exact integer arithmetic by an independent implementation (Appendix~\ref{app:definitions}); boundary maps are checked to satisfy $\partial^2=0$ in every computation. No numerical homology is used.
\end{remark}

\subsection{Monotonicity and plateaus}
\begin{lemma}[Lemma M: occurrence sets are monotone]\label{lem:M}
For every realised state $x$, $D(x)\subseteq D(C(x))$. In particular every orbit has a non-decreasing sequence $D_0\subseteq D_1\subseteq\cdots$ and every state of it is realised.
\end{lemma}
\prov{CW-0425}{A4-T3a}{PROVED}{FULL}
\begin{proof}
Let $G=bh$ as in \cref{lem:basic}, so $C(x)=(ah+L^2)/G$. For $d\in D(x)$ we have $d^2\mid a$ and $d\mid L$, so $d^2\mid ah+L^2$, and $\gcd(d,G)=1$ because $d\equiv1\pmod G$. Let $g=\gcd(ah+L^2,G)$; then $C(x)=a'/b'$ with $a'=(ah+L^2)/g$, $b'=G/g$. Since $g\mid G$ is prime to $d$, still $d^2\mid a'$; and $b'\mid G\mid d-1$ gives $d\equiv1\pmod{b'}$. Hence $d\in D(C(x))$. \qedhere
\end{proof}

\begin{definition}[events and plateaus]\label{def:event}
A step $x_n\to x_{n+1}$ is a \emph{$D$-event} if $D_n\ne D_{n+1}$; a maximal run of steps without events is a \emph{plateau} (or \emph{stage}). The sets $N=D_{n+1}\setminus D_n$ are the \emph{new occurrences}.
\end{definition}
Between events $L$ and $G$ are constant, so a plateau is an arithmetic progression $x_{n+j}=x_n+jL^2/G$; the sieve of \cref{sec:events} reads the length of the plateau from one congruence.
